\documentclass[11pt]{article}
\usepackage{qfffcalc}

\QFFFMetadata
  {$\mathrm{SU}(3)+5F$}
  {\frac12}
  {$\mathrm{SU}(5)\times\mathrm{U}(1)_B$}
  {$\mathrm{SU}(5)\times\mathrm{SU}(2)\times\mathrm{U}(1)_q$}
  {$O(t^{11})$}
  {Hilbert series and plethystic logarithm}
  {Complete}

\begin{document}
\MakeQFFFTitle

\section{Conventions}
\[
[a_1,a_2,a_3,a_4;b]
=[a_1,a_2,a_3,a_4]_{\mathrm{SU}(5)}\otimes[b]_{\mathrm{SU}(2)},
\qquad t^{2R}.
\]
The fugacity $q$ denotes the raw ultraviolet abelian charge.  Non-abelian unrefinement replaces each character by its dimension.

\section{Highest-weight generating function}
\begin{align}
\mathrm{HWG}
={}&\PE\!\left[(\mu_1\mu_4+\nu^2+1)t^2
+\nu(q\mu_2+q^{-1}\mu_3)t^3
+\mu_2\mu_3t^4-\nu^2\mu_2\mu_3t^6\right],
\label{eq:hwg-pe}
\end{align}
\begin{align}
\mathrm{HWG}
={}&\frac{1-\nu^2\mu_2\mu_3t^6}
{(1-t^2)(1-\nu^2t^2)(1-\mu_1\mu_4t^2)}
\nonumber\\[-1mm]
&\times\frac{1}
{(1-\nu q\mu_2t^3)(1-\nu q^{-1}\mu_3t^3)(1-\mu_2\mu_3t^4)}.
\label{eq:hwg-rational}
\end{align}

\section{Highest-weight expansion}
\begin{align*}
\mathrm{HWG}(t,q;\mu,\nu)={}&1
+\bigl(1+\nu^2+\mu_1\mu_4\bigr)t^2
+\bigl(q^{-1}\mu_3\nu+q\mu_2\nu\bigr)t^3\\
&+\Bigl(1+\nu^2+\nu^4+\mu_2\mu_3+\mu_1\mu_4
 +\mu_1\mu_4\nu^2+\mu_1^2\mu_4^2\Bigr)t^4\\
&+\Bigl(q^{-1}\mu_3\nu+q^{-1}\mu_3\nu^3
 +q\mu_2\nu+q\mu_2\nu^3\\
&\hspace{1.4cm}+q^{-1}\mu_1\mu_3\mu_4\nu
 +q\mu_1\mu_2\mu_4\nu\Bigr)t^5\\
&+\Bigl(1+\nu^2+\nu^4+\nu^6+q^{-2}\mu_3^2\nu^2
 +\mu_2\mu_3+\mu_2\mu_3\nu^2\\
&\hspace{1.4cm}+q^2\mu_2^2\nu^2+\mu_1\mu_4
 +\mu_1\mu_4\nu^2+\mu_1\mu_4\nu^4\\
&\hspace{1.4cm}+\mu_1\mu_2\mu_3\mu_4+\mu_1^2\mu_4^2
 +\mu_1^2\mu_4^2\nu^2+\mu_1^3\mu_4^3\Bigr)t^6
+O(t^7).
\end{align*}
The complete expansion through $t^{10}$ is given in Appendix~A.

\section{Hilbert series}
\subsection{Character-valued series}
\begin{align*}
H(t,q)={}&[0,0,0,0;0]\\
&+\bigl([0,0,0,0;0]+[0,0,0,0;2]+[1,0,0,1;0]\bigr)t^2\\
&+\bigl(q^{-1}[0,0,1,0;1]+q[0,1,0,0;1]\bigr)t^3\\
&+\Bigl([0,0,0,0;0]+[0,0,0,0;2]+[0,0,0,0;4]
+[0,1,1,0;0]\\
&\hspace{1.2cm}+[1,0,0,1;0]+[1,0,0,1;2]+[2,0,0,2;0]\Bigr)t^4\\
&+\Bigl(q^{-1}[0,0,1,0;1]+q^{-1}[0,0,1,0;3]
+q^{-1}[1,0,1,1;1]\\
&\hspace{1.2cm}+q[0,1,0,0;1]+q[0,1,0,0;3]
+q[1,1,0,1;1]\Bigr)t^5\\
&+\Bigl(q^{-2}[0,0,2,0;2]+[0,0,0,0;0]+[0,0,0,0;2]
+[0,0,0,0;4]\\
&\hspace{1.2cm}+[0,0,0,0;6]+[0,1,1,0;0]+[0,1,1,0;2]
+[1,0,0,1;0]\\
&\hspace{1.2cm}+[1,0,0,1;2]+[1,0,0,1;4]+[1,1,1,1;0]
+[2,0,0,2;0]\\
&\hspace{1.2cm}+[2,0,0,2;2]+[3,0,0,3;0]+q^2[0,2,0,0;2]\Bigr)t^6
+O(t^7).
\end{align*}

\subsection{$q$-refined dimension series}
\input{data/f5k12_hdim_full.tex}

\subsection{Unrefined series}
\begin{align}
H(t)={}&1+28t^2+40t^3+380t^4+820t^5+3656t^6+8460t^7
\nonumber\\
&+27076t^8+60620t^9+160068t^{10}+O(t^{11}).
\label{eq:unrefined-hs}
\end{align}

\paragraph{Literature comparison.}
The coefficients through $t^6$ agree with equation (12.5) of Ferlito--Hanany--Mekareeya--Zafrir for the $n=3$ member of the $\mathrm{SU}(n)_{\pm1/2}$, $N_f=2n-1$ sequence.  The terms from $t^7$ through $t^{10}$ extend the displayed order.

\section{Plethystic logarithm}
\subsection{Character-valued series}
\begin{align*}
\PL[H(t,q)]={}&
\bigl([0,0,0,0;0]+[0,0,0,0;2]+[1,0,0,1;0]\bigr)t^2\\
&+\bigl(q^{-1}[0,0,1,0;1]+q[0,1,0,0;1]\bigr)t^3\\
&-\bigl(2[0,0,0,0;0]+[1,0,0,1;0]\bigr)t^4\\
&-\Bigl(q^{-1}[0,0,0,2;1]+2q^{-1}[0,0,1,0;1]
+q^{-1}[1,1,0,0;1]\\
&\hspace{1.2cm}+q[0,0,1,1;1]+2q[0,1,0,0;1]
+q[2,0,0,0;1]\Bigr)t^5\\
&-\Bigl(q^{-2}[0,1,0,1;0]+q^{-2}[1,0,0,0;2]
+[0,0,0,0;0]\\
&\hspace{1.2cm}+[0,0,0,0;2]+[0,1,1,0;0]+[0,1,1,0;2]
+[1,0,0,1;2]\\
&\hspace{1.2cm}+q^2[0,0,0,1;2]+q^2[1,0,1,0;0]\Bigr)t^6
+O(t^7).
\end{align*}

\subsection{Dimension-refined and unrefined series}
\input{data/f5k12_pl_dim_full.tex}
\input{data/f5k12_pl_unrefined_full.tex}

\begin{thebibliography}{9}
\bibitem{Ferlito2017}
G.~Ferlito, A.~Hanany, N.~Mekareeya and G.~Zafrir,
``3d Coulomb branch and 5d Higgs branch at infinite coupling,''
\emph{JHEP} \textbf{07} (2018) 061,
\href{https://arxiv.org/abs/1712.06604}{arXiv:1712.06604}.
\end{thebibliography}

\BeginQFFFAppendices

\clearpage
\begin{landscape}
\section{Complete highest-weight expansion through $t^{10}$}
{\small\input{data/f5k12_highest_weight_full.tex}}
\end{landscape}

\clearpage
\begin{landscape}
\section{Complete character Hilbert series through $t^{10}$}
{\small\input{data/f5k12_character_hs_full.tex}}
\end{landscape}

\clearpage
\begin{landscape}
\section{Complete character-valued plethystic logarithm through $t^{10}$}
{\footnotesize\input{data/f5k12_pl_character_full.tex}}
\end{landscape}

\end{document}
